\documentclass[12pt]{article}

\usepackage[utf8]{inputenc}
\usepackage[T1]{fontenc}
\usepackage{geometry}
\geometry{margin=1in}
\usepackage{graphicx}
\usepackage{amsmath,amssymb}
\usepackage{hyperref}
\hypersetup{colorlinks=true, linkcolor=blue, citecolor=blue, urlcolor=blue}
\usepackage{authblk}
\usepackage{abstract}
\usepackage{booktabs}
\usepackage{caption}
\usepackage[section]{placeins}

% NOTE: This manuscript is prepared with the standard LaTeX article class so
% that it compiles without external downloads. For submission to an AAS
% journal (RNAAS, PASP, ApJ, etc.) the \documentclass line and a small set
% of front-matter commands should be swapped for the aastex631 class; the
% body content (sections, equations, figures, references) requires no
% changes.

\title{\textbf{Uncertainty-Aware Colorization of Astronomical Images:\\
A Monte Carlo Dropout Approach with Perceptual Validation}}

\author[1]{Ali Teymouri}
\affil[1]{Independent Researcher}

\date{\today}

\begin{document}
\maketitle

\begin{abstract}
Composite color images of narrowband astronomical observations (the
``Hubble Palette,'' mapping SII, H$\alpha$, and OIII filters to red, green,
and blue channels) are produced by deterministic, hand-tuned stretches that
give no indication of per-pixel confidence. We present a pipeline that
instead learns this mapping with a compact convolutional neural network
trained under simulated channel dropout and sensor noise, and that uses
Monte Carlo Dropout at inference time to obtain, in addition to a color
estimate, a pixel-wise uncertainty map. We apply the method to Hubble Space
Telescope WFC3 narrowband observations of M16 (the Eagle Nebula / Pillars of
Creation), retrieved from the Mikulski Archive for Space Telescopes (MAST).
Qualitatively, the resulting uncertainty map is elevated at saturated stars,
frame edges, and cosmic-ray-contaminated pixels, and low across the
well-observed nebular body. We further describe a web-based perceptual
survey designed to test whether human raters judge model uncertainty
consistently with subjective color trustworthiness; we report preliminary
results from an initial, small rater pool and outline the data collection
still required before a statistically meaningful correlation can be claimed.
All code is released publicly to support reproduction and extension to
other targets.
\end{abstract}

\noindent\textbf{Keywords:} astronomical image processing --- deep learning ---
uncertainty quantification --- Bayesian neural networks --- Hubble Space Telescope

\section{Introduction}

Public-facing astronomical images are almost always composites: multiple
monochrome exposures, typically in narrowband filters isolating specific
emission lines (e.g., [S~II] $\lambda6716,6731$, H$\alpha$ $\lambda6563$,
[O~III] $\lambda5007$), are individually stretched and assigned to color
channels. The de facto standard for emission nebulae is the ``Hubble
Palette'' (SII$\to$R, H$\alpha\to$G, OIII$\to$B), popularized by the Hubble
Heritage Project. While visually effective, this process is deterministic
and offers no mechanism to communicate where the resulting color is
well-constrained by the data and where it is not --- for example, in regions
where one band is missing, saturated, or contaminated by cosmic rays.

Bayesian deep learning offers a practical route to per-pixel confidence
estimates without the computational cost of a full Bayesian neural network.
\citet{gal2016dropout} showed that a network trained with Dropout can be
interpreted as an approximation to a deep Gaussian process, and that keeping
Dropout \emph{active} at inference time and averaging multiple stochastic
forward passes (Monte Carlo Dropout, MC-Dropout) yields both a predictive
mean and a predictive variance, the latter serving as an uncertainty
estimate. MC-Dropout has been applied broadly in medical imaging and
autonomous driving, where it is important to know not just what a model
predicts but how much to trust that prediction; its application to
astronomical image colorization, to our knowledge, is not established in
the mainstream literature, motivating the exploratory pipeline described
here.

This paper makes three contributions: (1) an end-to-end, reproducible
pipeline for MC-Dropout-based uncertainty-aware colorization of multi-band
astronomical imagery, from raw MAST retrieval through registration,
training, and tiled inference; (2) a qualitative analysis of the resulting
uncertainty map on a real HST dataset of M16; and (3) the design (and
preliminary results) of a perceptual survey intended to test whether the
model's uncertainty is perceptually meaningful to human observers.

\section{Data}

We use narrowband imaging of M16 (NGC 6611, the Eagle Nebula) obtained with
the Wide Field Camera 3 (WFC3/UVIS) aboard the Hubble Space Telescope,
retrieved via \texttt{astroquery.mast} \citep{astroquery, astropy2018}. Three
filters were used, each isolating a different ionic emission line:

\begin{itemize}
  \item \textbf{F502N} (\textrm{[O~III]}, $\lambda5007$)
  \item \textbf{F656N} (H$\alpha$, $\lambda6563$)
  \item \textbf{F673N} (\textrm{[S~II]}, $\lambda6716,6731$)
\end{itemize}

The OIII and SII exposures originate from the same observing proposal
(proposal ID 13926) and share a common pointing. No H$\alpha$ exposure from
the same proposal was available in the archive; the closest-coverage
H$\alpha$ observation (proposal 16359) was retrieved instead. Because its
footprint does not fully coincide with the OIII/SII pointing, the three
bands were reprojected onto a common World Coordinate System (WCS) grid
using \texttt{reproject.reproject\_interp} \citep{reproject}, and the image
was cropped to the largest axis-aligned rectangle with 100\% valid coverage
in all three bands, found via a maximal-rectangle-in-binary-matrix search
over the union-validity mask. This crop is necessary because the naive
bounding box of the valid-pixel mask still contains invalid corners when the
input footprints are mutually rotated.

\section{Methods}

\subsection{Baseline colorization}

As a reference, we construct a deterministic Hubble Palette composite: each
registered, cropped band is independently stretched with a 1st--99.5th
percentile linear stretch and assigned to its corresponding RGB channel.
This baseline has no notion of confidence and serves only as a visual and
numerical reference point for the learned model.

\subsection{MC-Dropout network}

We use a compact fully convolutional network (four convolutional layers,
32--64 channels, ReLU activations, 1$\times$1 output convolution with a
sigmoid nonlinearity) with \texttt{Dropout2d} ($p=0.3$) inserted after each
hidden layer. Because no independently verified ``true color'' exists for
these data, we cannot train a supervised colorization model in the
conventional sense. Instead we frame the learning problem as
\emph{reconstruction under simulated data degradation}: during training,
each input patch has, with probability 0.4, one of its three bands set to
zero (simulating a missing/uncovered band, as encountered with the real
H$\alpha$ footprint mismatch), plus additive Gaussian noise
($\sigma=0.05$) applied to all channels. The training target is the clean,
undegraded three-band patch. This framing gives the resulting uncertainty a
concrete operational meaning: it should be elevated wherever the network is
extrapolating from incomplete or noisy information, which is directly
relevant to the missing-coverage problem observed in Section 2.

The network is trained with Adam (learning rate $10^{-3}$) and an MSE loss
on $128\times128$ patches randomly sampled from the registered, percentile
stretched bands, for 30 epochs (4000 patches/epoch, batch size 16), on a
single NVIDIA RTX 4060 laptop GPU.

\subsection{Uncertainty inference}

At inference, Dropout is kept active (\texttt{model.train()} mode) and the
full image is processed in overlapping $256\times256$ tiles (stride 192) to
avoid seam artifacts. Each tile is passed through the network $N=30$ times;
the per-pixel mean across passes is taken as the color estimate and the
per-pixel standard deviation as the uncertainty estimate. Tiles are
recombined by averaging in overlap regions.

\subsection{Perceptual validation survey}

To assess whether the resulting uncertainty map corresponds to human
judgments of color trustworthiness, we selected 14 $180\times180$-pixel
patches spanning the observed range of mean per-patch uncertainty (from the
lowest to the highest decile), with display order randomized per session.
Human raters, blind to the underlying uncertainty values, scored each patch
on a 5-point scale (1 = ``completely artificial,'' 5 = ``completely natural
and trustworthy''). The instrument is a self-contained, browser-based
survey (Supplementary material; \texttt{survey/uncertainty\_survey.html})
that requires no server-side infrastructure. We planned to correlate the
mean human rating per patch against the model's mean uncertainty for that
patch using Spearman's rank correlation coefficient $\rho$.

\section{Results}

\subsection{Registration and baseline}

After reprojection, the OIII and SII bands (originating from the same
pointing) retained 93\% pixel coverage relative to the reference frame,
while the independently-pointed H$\alpha$ band retained only 26--33\%
coverage depending on the specific archival exposure tested. The
maximal-rectangle crop yields a final working image of
$1920\times3873$ pixels with 100\% valid coverage in all three bands.

\subsection{Uncertainty map}

Qualitatively, the MC-Dropout uncertainty map (Figure~1, not reproduced
here; see \texttt{survey/} outputs) shows three consistent patterns: (i)
elevated uncertainty at a saturated foreground star, (ii) elevated
uncertainty along image borders, where the network has reduced spatial
context, and (iii) elevated uncertainty at scattered single-pixel and
short-streak features consistent with cosmic-ray hits, which by
construction appear in only one of the three independently-exposed bands
and therefore constitute exactly the kind of ``missing/inconsistent
information'' the training procedure was designed to flag. The bulk of the
nebular body, where all three bands are well exposed and mutually
consistent, shows low uncertainty.

\subsection{Preliminary perceptual survey results}

At the time of writing, $n=3$ raters had completed the survey (partial
completion permitted). The Spearman correlation between mean human rating
and model uncertainty across the 14 patches was $\rho \approx 0.00$
($p \approx 1.0$), i.e., no significant relationship was detected. Given
$n=3$, this is expected to be underpowered and should not be interpreted as
evidence against the method; a minimum of 10--15 raters is planned before
drawing conclusions. One qualitative observation is noteworthy: the single
patch with the highest model uncertainty (containing the saturated star)
also showed the largest disagreement \emph{among raters} (ratings of 1, 4,
and 4 on the 5-point scale), suggesting that this patch was subjectively
ambiguous to humans as well, even though the average rating alone does not
yet show a monotonic trend with uncertainty. Data collection is ongoing;
this section will be updated with a full statistical analysis
(correlation, inter-rater agreement, and a scatter plot of rating vs.\
uncertainty) once a sufficient sample is collected.

\section{Discussion}

Several limitations should be noted. First, because no ground-truth color
exists, the network's training target is itself a deterministic stretch of
the observed bands rather than an independently verified ``correct''
color; the uncertainty we estimate is therefore best understood as
\emph{data-completeness/consistency uncertainty} (how much the network
relies on extrapolation from degraded input) rather than a calibrated
statistical confidence interval on true physical flux. Second, the
perceptual survey, as of this draft, has too few respondents for a
statistically meaningful correlation; the survey is designed to scale to
additional raters without further engineering (see
\texttt{07\_generate\_survey\_patches.py} and the accompanying HTML
instrument), and this is the immediate next step in the project. Third,
results are demonstrated on a single target (M16); generalization to other
nebulae, filter sets, and instruments (e.g., JWST NIRCam) remains to be
tested.

\section{Conclusion and Future Work}

We have demonstrated a complete, reproducible pipeline for MC-Dropout-based
uncertainty-aware colorization of multi-band Hubble Space Telescope
imagery, including automated data retrieval, robust registration under
mismatched fields of view, a degradation-aware training scheme that gives
the resulting uncertainty a concrete operational interpretation, and a
lightweight perceptual survey instrument for human validation. Immediate
next steps are (1) completing perceptual data collection to a statistically
adequate sample size, (2) extending the pipeline to additional targets and
instruments, and (3) exploring alternative uncertainty-estimation
techniques (deep ensembles, evidential deep learning) as a comparison point
for MC-Dropout. All code, the survey instrument, and this manuscript source
are released publicly (see repository).

\section*{Acknowledgments}
Based on observations made with the NASA/ESA Hubble Space Telescope,
obtained from the data archive at the Space Telescope Science Institute
(STScI), which is operated by the Association of Universities for Research
in Astronomy, Inc., under NASA contract NAS 5-26555. This research made use
of \texttt{astropy} \citep{astropy2018}, \texttt{astroquery}
\citep{astroquery}, and \texttt{reproject} \citep{reproject}.

\begin{thebibliography}{9}

\bibitem[Gal \& Ghahramani(2016)]{gal2016dropout}
Gal, Y., \& Ghahramani, Z. 2016, in Proceedings of the 33rd International
Conference on Machine Learning (ICML), \textit{Dropout as a Bayesian
Approximation: Representing Model Uncertainty in Deep Learning}

\bibitem[Astropy Collaboration(2018)]{astropy2018}
Astropy Collaboration, Price-Whelan, A. M., Sip\H{o}cz, B. M., et al.\ 2018,
\textit{AJ}, 156, 123

\bibitem[Ginsburg et al.(2019)]{astroquery}
Ginsburg, A., Sip\H{o}cz, B. M., Brasseur, C. E., et al.\ 2019,
\textit{AJ}, 157, 98 (\texttt{astroquery})

\bibitem[Robitaille(2019)]{reproject}
Robitaille, T. 2019, \textit{reproject: Astronomical image reprojection in
Python}, Astropy Project

\end{thebibliography}

\end{document}
